\begin{abstract}
本文提出原型路由图提示方法 \method，利用支持集原型将可学习提示分配给不同节点。在每类仅有少量标注的条件下，任务适配不仅需要学习提示内容，还需要确定提示应当作用于哪些节点。\method 从聚合节点属性构造类别参照，在保留支持集锚点的同时利用未标注图特征进行校正，并将节点与原型的关联转换为固定路由矩阵。该矩阵组合可学习类别提示，将其注入冻结图卷积编码器，最终采用原型分类。我们在 Cora、CiteSeer、PubMed 上分别评估 1、3、5-shot 设置，对各条件使用配对的支持抽样与训练种子。5-shot 下，三个数据集的 Accuracy 分别为 73.02\%、66.17\%、70.98\%；相对同配方共享提示，原型路由在对应条件下分别带来 36.94、40.45、3.02 个百分点的提升。这些逐条件配对结果支持原型关系对提示分配的组织价值，结论限于所评估的传导式节点分类设置。
\end{abstract}

\section{引言}
预训练图神经网络为标注稀缺条件下的节点分类提供了已有表示。图提示通过少量任务参数完成适配，具体形式包括特征修改、提示结构和任务读出~\citep{gppt2022,gpf2023,allinone2023,graphprompt2023}。除选择注入位置之外，适配方法还需要回答：提示的贡献应当如何分配给不同节点？同一图中的节点在下游分类任务中具有不同的角色，因此提示分配本身就是值得研究的设计问题。

共享提示表达全图一致的调整。更灵活的方法可以学习输入相关的组合、在子图层面分配提示，或通过策略选择被提示的节点~\citep{gpf2023,supt2026,relief2025}。这些工作说明，节点相关的分配本身并非新概念。本文关注一个更具体的问题：\emph{任务已有的少量标注样本，能否直接定义有效的路由系数，使优化集中于提示内容？} 支持样本提供了当前任务的类别参照，这些参照可以在适配之前构造，无需额外训练路由网络。

基于这一问题，本文提出 \method，以支持集锚定的原型作为提示分配的任务坐标。类别关联相近的节点接收相关的类别提示组合。方法首先从聚合图属性构造路由表示，以标注支持节点初始化原型，并在保留支持锚点的条件下利用未标注特征校正原型。节点与原型的关联形成固定路由矩阵，在冻结 GCN 内部组合可学习提示向量。训练仅优化提示内容，最终采用原型分类。图~\ref{fig:architecture} 概括了固定分配与内容学习的分工。

实验围绕两个问题组织。首先，完整适配在本地配对运行中相对冻结预训练和全量微调表现如何，其 Accuracy 在文献报告的提示结果中处于什么位置？其次，在注入位置和训练配方固定时，原型路由能否提供超出共享提示的价值？本地运行覆盖 Cora、CiteSeer、PubMed 的 1、3、5-shot；外部 Accuracy 取自 \citet{mtg2026}，保留其原始协议。每项结果均对应独立的数据集和 shot 数。相对冻结预训练，PubMed 上的收益明显小于 Cora 与 CiteSeer。

本文的主要贡献如下：
\begin{itemize}
\item 提出原型路由图提示构造，由支持集锚定的类别关系在提示优化前确定节点相关的分配方式。
\item 给出包含 $C\times128$ 个可学习提示参数的隐藏层实现，无需可学习路由网络。
\item 提供配对适配比较与共享提示消融，逐条件报告 seed 波动及配对不确定性。
\end{itemize}
贡献聚焦原型关系对隐藏层提示的组织作用，不将原型或节点相关提示视为新的基本构件。

\section{相关工作}
\paragraph{图表示与预训练。}
图卷积、邻域采样与图注意力提供了不同的节点表示构造方式~\citep{gcn2017,sage2017,gat2018}。图自监督目标包括局部与全局信息最大化~\citep{dgi2019}、基于增广的对比学习~\citep{graphcl2020}以及掩码重建~\citep{graphmae2022}。预训练目标会影响下游迁移，联合节点级与图级目标也已有专门研究~\citep{pretrain2020}。本文固定每个图的预训练权重，研究下游适配，不提出新的预训练目标。

\paragraph{预训练与下游任务对齐。}
GPPT 与 GraphPrompt 通过任务重构和提示连接预训练与下游预测~\citep{gppt2022,graphprompt2023}。All in One 通过提示向量、结构和插入规则处理多层级任务~\citep{allinone2023}；MultiGPrompt 结合多个预训练任务及下游提示~\citep{multigprompt2024}。HGPrompt 与 HetGPT 将提示学习拓展到异构图设置~\citep{hgprompt2024,hetgpt2024}。这些工作关注任务或图类型之间的对齐，本文则在少样本节点分类中研究更具体的问题：如何利用任务类别参照在节点间分配提示。

\paragraph{自适应提示分配。}
GPF-plus 学习输入相关的提示组合~\citep{gpf2023}，SUPT 在子图层面分配提示~\citep{supt2026}，RELIEF 联合学习节点选择与提示内容~\citep{relief2025}。ProNoG 面向非同配图研究条件化的节点专属提示~\citep{pronog2025}，DAGPrompT 关注分布感知的图适配~\citep{dagprompt2025}。属性引导的动态提示也已用于提升模型应对图扰动的鲁棒性~\citep{adprompt2026}。与学习分配函数的思路不同，\method 预先计算支持集条件下的原型关系，并在提示训练期间固定路由。上述工作用于界定设计空间，并非全部纳入本文的实验比较。

\paragraph{类别原型与评估。}
Prototypical Networks 利用支持集类别代表进行预测~\citep{proto2017}，半监督扩展则将未标注样本用于原型构造~\citep{semiproto2018}。本文利用这些已有思想构造路由的任务参照：路由原型位于聚合属性空间，预测原型来自适配后的编码器输出。ProG 强调评估协议的明确性~\citep{prog2024}。\citet{mtg2026} 报告了不同预训练策略下的适配结果；本文引用其基线 Accuracy 作为文献参照，同时保留本地配对消融以检验路由价值。

\section{问题定义}
给定属性图 $G=(V,E,X)$，节点数为 $n$，特征矩阵为 $X\in\mathbb R^{n\times d}$，任务类别数为 $C$。每个类别 $c$ 的支持集 $S_c$ 包含 $k$ 个标注节点，其中 $k\in\{1,3,5\}$。记 $S=\bigcup_c S_c$，查询集 $T$ 与 $S$ 不相交。完整图结构与所有节点特征可见，拟合和超参数选择仅接收 $S$ 内的标签；查询标签在预测封存后用于评分。

编码器采用预训练两层 GCN~\citep{gcn2017}，隐藏维度与输出维度均为 128，在 \method 中保持冻结。本文属于同图、已知类别的传导式分类：查询特征参与图聚合、原型校正和全图中心化，不属于归纳式未见类别评估。文中的一个\emph{实验条件}固定数据集、shot 数、方法及预训练权重。

\section{原型路由图提示}
\label{sec:method}
\begin{figure*}[t]
\centering
\fbox{\begin{minipage}[c][32mm][c]{0.92\linewidth}
\centering
\lang{Architecture figure placeholder}{主架构图占位符}\\[3mm]
\lang{Graph attributes and support labels $\longrightarrow$ anchored prototypes $\longrightarrow$ fixed routing}{图属性与支持标签 $\longrightarrow$ 支持集锚定原型 $\longrightarrow$ 固定路由}\\[3mm]
\lang{Fixed routing + trainable prompt tokens $\longrightarrow$ frozen GCN $\longrightarrow$ prototype prediction}{固定路由与可学习提示 $\longrightarrow$ 冻结 GCN $\longrightarrow$ 原型分类}
\end{minipage}}
\caption{\lang{Overview of \method. Support labels initialize class references in aggregated attribute space. Unlabeled graph features refine these references, and the resulting node--prototype associations determine a fixed routing matrix. Only the class prompt tokens are optimized; the routed prompts enter the first hidden layer of a frozen GCN. Query labels are excluded from fitting and selection. The final architecture artwork is intentionally left as a placeholder.}{\method 总体流程。支持标签在聚合属性空间中初始化类别参照，未标注图特征用于校正参照，节点与原型的关联决定固定路由矩阵。仅类别提示向量参与优化，路由后的提示注入冻结 GCN 的第一隐藏层。拟合与选参不使用查询标签。主架构图按要求暂用占位符。}}
\label{fig:architecture}
\end{figure*}

\paragraph{方法概览。}
方法将分配与内容分开处理：首先，在聚合属性空间中用支持集锚定原型定义类别关系；其次，将这些关系转换为固定路由矩阵，为每个节点组合提示；最后，利用支持监督在冻结编码器内学习提示向量。以下各节依照这一计算顺序展开。
\subsection{支持集锚定的类别参照}
令 $\widehat A=\widetilde D^{-1/2}(A+I)\widetilde D^{-1/2}$ 为带自环的归一化邻接算子。由两跳聚合属性构造路由表示：
\begin{equation}
\begin{aligned}
X^{(2)}&=\widehat A^2X,\\
R&=\rownorm\!\left(X^{(2)}-\mathbf1\,\mean_{v\in V}X^{(2)}_v\right).
\end{aligned}
\label{eq:attributes}
\end{equation}
其中 $\rownorm$ 对每一行进行欧氏归一化。路由使用图属性及结构，不依赖额外学习的网络。类别初始参照为
\begin{equation}
\mu_c^{(0)}=\frac{1}{k}\sum_{i\in S_c}R_i.
\label{eq:anchor}
\end{equation}

为利用未标注图信息并保留标注参照，记 $U_V=V\setminus S$。从 $\mu^{(0)}$ 开始，令 $i\in U_V$、$t=0,1,2$，进行三轮固定校正：
\begin{align}
a_{ic}^{(t)}&=\frac{\exp(R_i^\top\norm{\mu_c^{(t)}}/0.05)}{\sum_{c'}\exp(R_i^\top\norm{\mu_{c'}^{(t)}}/0.05)},\label{eq:assignment}\\
\mu_c^{(t+1)}&=0.5\mu_c^{(0)}+0.5\frac{\sum_{i\in U_V}a_{ic}^{(t)}R_i}{\max(\sum_{i\in U_V}a_{ic}^{(t)},10^{-8})}.\label{eq:refinement}
\end{align}
每轮锚点始终为最初的支持原型，而不是上一轮原型。这是固定的原型构造过程，本文不赋予其期望最大化目标或收敛保证。计算不使用 $U_V$ 中节点的标签。

\subsection{从类别关系到固定路由}
以校正后的原型计算温度为 $\tau_r=0.1$ 的节点—类别关联：
\begin{equation}
\begin{aligned}
Q_{ic}&=\softmax_c\!\left(R_i^\top\norm{\mu_c^{(3)}}/\tau_r\right),\\
B&=\widehat A\left(Q-\mathbf1\,\mean_{v\in V}Q_v\right).
\end{aligned}
\label{eq:routing}
\end{equation}
中心化表示节点类别关联相对全图平均关联的偏离，邻域聚合进一步整合相邻节点的关联模式。$Q$ 是按行归一化的关联矩阵，$B$ 则包含可为负的提示组合系数，不能解释为最终类别概率。每个支持集仅计算一次 $B$，提示训练期间保持不变。

\subsection{冻结隐藏计算中的提示适配}
令 $H^{\mathrm{pre}}$ 为原图经过 GCN 第一层得到的预激活表示，学习零初始化的矩阵 $P\in\mathbb R^{C\times128}$。记 $s$ 为 $H^{\mathrm{pre}}$ 的逐通道样本标准差，下界设为 $10^{-6}$，则注入后的隐藏表示为
\begin{equation}
\widetilde H^{\mathrm{pre}}=H^{\mathrm{pre}}+10\,[B\tanh(P)]\odot s.
\label{eq:injection}
\end{equation}
尺度向量按节点广播。经过原有 LeakyReLU 与第二层图卷积后得到 $Z$。支持监督优化提示内容，路由与编码器参数保持固定；零初始化使训练起点恢复冻结编码器的原始计算。Cora、CiteSeer、PubMed 分别包含 896、768、384 个可学习提示参数。

\subsection{训练与预测}
使用 Adam 优化提示矩阵 100 步，学习率为 0.01，权重衰减为 0：
\begin{equation}
\begin{aligned}
\mathcal L(P)={}&\frac{1}{|S|}\sum_{i\in S}\operatorname{CE}(\ell_i(P),y_i)\\
&+10^{-3}\mean(P\odot P).
\end{aligned}
\label{eq:loss}
\end{equation}
支持 logits $\ell_i$ 采用余弦原型分类，路由消融始终使用相同的原型读出。支持集留一构造、预测归一化及温度等复现细节统一放在附录~\ref{app:protocol}。

路由只需预先计算一次。在稀疏图上，两次属性聚合的复杂度为 $O(|E|d)$，每轮原型校正为 $O(nCd)$，每步提示组合 $B\tanh(P)$ 为 $O(nC\cdot128)$，此外还包含冻结编码器的前向与反向计算。冻结权重减少可学习参数，但不会消除经过第二层卷积的梯度计算。因此本文报告参数量，不声称未经测量的运行时间优势。

\section{实验}
\subsection{协议与比较方法}
\label{sec:protocol}
\begin{table*}[t]
\centering
\linespread{1}\selectfont
\captionsetup{font=small,skip=5pt}
\caption{\lang{\textbf{Accuracy on few-shot node classification.} Higher is better ($\uparrow$).}{\textbf{少样本节点分类的 Accuracy。}越高越好（$\uparrow$）。}}
\label{tab:main-accuracy}
\fontsize{9}{11}\selectfont
\setlength{\tabcolsep}{1pt}
\renewcommand{\arraystretch}{1.18}
\setlength{\heavyrulewidth}{0.8pt}
\setlength{\lightrulewidth}{0.4pt}
\begin{tabular*}{\linewidth}{@{\extracolsep{\fill}}l*{9}{r}@{}}
\toprule
\lang{Method}{方法} & \multicolumn{3}{c}{\textbf{Cora}} & \multicolumn{3}{c}{\textbf{CiteSeer}} & \multicolumn{3}{c}{\textbf{PubMed}} \\
\cmidrule(lr){2-4}\cmidrule(lr){5-7}\cmidrule(l){8-10}
 & 1-shot & 3-shot & 5-shot & 1-shot & 3-shot & 5-shot & 1-shot & 3-shot & 5-shot \\
\midrule
\multicolumn{10}{@{}l}{\textit{\lang{Literature-reported baselines}{文献报告基线}~\citep{mtg2026}}} \\[2pt]
Supervised & \meanstd{26.56}{5.55} & \meanstd{37.79}{9.16} & \meanstd{50.25}{8.37} & \meanstd{21.78}{7.32} & \meanstd{35.18}{6.86} & \meanstd{41.22}{6.30} & \meanstd{39.37}{16.34} & \meanstd{57.33}{4.64} & \meanstd{67.88}{2.18} \\
Fine-tuning & \meanstd{52.61}{1.73} & \meanstd{51.97}{2.84} & \meanstd{62.66}{3.55} & \meanstd{35.05}{4.37} & \meanstd{45.08}{2.09} & \meanstd{39.54}{3.54} & \meanstd{46.74}{14.89} & \meanstd{65.40}{3.00} & \bestmeanstd{70.91}{4.87} \\
\midrule
GPPT & \meanstd{43.15}{9.44} & \meanstd{43.84}{6.11} & \meanstd{51.98}{3.43} & \meanstd{37.26}{6.17} & \meanstd{42.34}{8.31} & \meanstd{45.77}{7.41} & \secondmeanstd{48.31}{17.72} & \meanstd{67.43}{2.96} & \meanstd{66.97}{3.70} \\
Gprompt & \bestmeanstd{56.66}{11.22} & \bestmeanstd{63.78}{5.77} & \bestmeanstd{69.03}{3.61} & \secondmeanstd{53.21}{10.94} & \secondmeanstd{60.00}{6.18} & \secondmeanstd{66.13}{1.64} & \meanstd{39.74}{15.35} & \meanstd{66.68}{3.53} & \meanstd{67.87}{2.08} \\
All-in-one & \meanstd{52.39}{10.17} & \meanstd{48.09}{4.83} & \meanstd{30.36}{13.48} & \meanstd{40.41}{2.80} & \meanstd{48.09}{8.18} & \meanstd{27.93}{10.59} & \meanstd{45.17}{6.45} & \meanstd{65.79}{5.79} & \meanstd{46.16}{15.83} \\
GPF & \meanstd{38.57}{5.41} & \meanstd{34.84}{19.83} & \meanstd{35.43}{1.02} & \meanstd{31.16}{8.05} & \meanstd{25.92}{12.30} & \meanstd{25.12}{3.01} & \bestmeanstd{49.99}{8.86} & \bestmeanstd{71.20}{2.82} & \meanstd{68.96}{3.99} \\
GPF-plus & \secondmeanstd{55.77}{10.30} & \secondmeanstd{56.38}{5.37} & \secondmeanstd{66.22}{6.20} & \bestmeanstd{59.67}{11.87} & \bestmeanstd{72.48}{5.63} & \bestmeanstd{75.73}{2.19} & \meanstd{46.64}{18.97} & \secondmeanstd{70.85}{4.03} & \secondmeanstd{69.59}{4.33} \\
\midrule
\multicolumn{10}{@{}l}{\textit{\lang{Local measurements (different protocol)}{本地测量（协议不同）}}} \\[2pt]
\frozenname & \meanstd{23.17}{5.48} & \meanstd{29.42}{3.47} & \meanstd{33.56}{5.03} & \meanstd{21.14}{2.87} & \meanstd{25.29}{2.16} & \meanstd{27.38}{2.84} & \secondmeanstd{55.44}{11.42} & \secondmeanstd{65.55}{8.11} & \secondmeanstd{68.37}{3.41} \\
\ftname & \secondmeanstd{45.41}{10.18} & \secondmeanstd{61.54}{6.78} & \secondmeanstd{67.50}{5.06} & \secondmeanstd{35.76}{6.80} & \secondmeanstd{42.22}{5.88} & \secondmeanstd{47.89}{4.85} & \meanstd{48.77}{8.54} & \meanstd{59.75}{7.97} & \meanstd{63.96}{5.18} \\
\midrule\rowcolor{black!5}
\textbf{\method} & \bestmeanstd{54.59}{9.18} & \bestmeanstd{70.76}{2.76} & \bestmeanstd{73.02}{1.95} & \bestmeanstd{47.12}{6.31} & \bestmeanstd{61.20}{4.98} & \bestmeanstd{66.17}{2.16} & \bestmeanstd{58.70}{9.78} & \bestmeanstd{68.66}{4.58} & \bestmeanstd{70.98}{2.46} \\
\bottomrule
\end{tabular*}
\par\vspace{3pt}
\begin{minipage}{\linewidth}\fontsize{8}{9.5}\selectfont\raggedright \lang{Mean $\pm$ std. (\%). Literature panel: \citet{mtg2026}, Tables 4--6, five trials, 90\% test allocation, best reported across pretraining strategies. Local panel: fixed edge-pretrained checkpoint and official test split; three training seeds averaged per support draw, then mean/sample std. over 20 draws. \textbf{Bold}/\underline{underline}: best/second within each panel and column only. Protocols differ; no cross-panel ranking, pairing, or pooled average.}{均值 $\pm$ 标准差（\%）。文献面板取自 \citet{mtg2026} 表 4--6：5 次抽样、90\% 测试集、不同预训练策略中的最佳报告结果。本地面板使用固定边预测预训练权重与官方测试划分，先平均每个支持抽样的 3 个训练 seed，再统计 20 次抽样的均值／样本标准差。\textbf{粗体}／\underline{下划线}仅表示同面板同列最优／次优；不跨协议排名、配对检验或求平均。}\end{minipage}
\end{table*}

\begin{table*}[t]
\centering
\linespread{1}\selectfont
\captionsetup{font=small,skip=5pt}
\caption{\lang{\textbf{Macro-F1 on few-shot node classification.} Higher is better ($\uparrow$).}{\textbf{少样本节点分类的 Macro-F1。}越高越好（$\uparrow$）。}}
\label{tab:main-macro_f1}
\fontsize{9}{11}\selectfont
\setlength{\tabcolsep}{1pt}
\renewcommand{\arraystretch}{1.18}
\setlength{\heavyrulewidth}{0.8pt}
\setlength{\lightrulewidth}{0.4pt}
\begin{tabular*}{\linewidth}{@{\extracolsep{\fill}}l*{9}{r}@{}}
\toprule
\lang{Method}{方法} & \multicolumn{3}{c}{\textbf{Cora}} & \multicolumn{3}{c}{\textbf{CiteSeer}} & \multicolumn{3}{c}{\textbf{PubMed}} \\
\cmidrule(lr){2-4}\cmidrule(lr){5-7}\cmidrule(l){8-10}
 & 1-shot & 3-shot & 5-shot & 1-shot & 3-shot & 5-shot & 1-shot & 3-shot & 5-shot \\
\midrule
\frozenname & \meanstd{22.51}{4.82} & \meanstd{29.07}{3.47} & \meanstd{33.07}{4.47} & \meanstd{20.39}{2.62} & \meanstd{24.32}{1.87} & \meanstd{26.45}{2.71} & \secondmeanstd{53.44}{10.49} & \secondmeanstd{63.19}{8.60} & \secondmeanstd{66.27}{3.55} \\
\ftname & \secondmeanstd{44.38}{10.07} & \secondmeanstd{60.94}{6.41} & \secondmeanstd{66.99}{4.52} & \secondmeanstd{30.99}{6.30} & \secondmeanstd{38.43}{5.80} & \secondmeanstd{44.62}{4.98} & \meanstd{44.52}{9.85} & \meanstd{57.64}{9.26} & \meanstd{62.53}{6.05} \\
\midrule
GPF & \meanstd{7.08}{2.07} & \meanstd{21.14}{6.64} & \meanstd{25.87}{6.19} & \meanstd{6.78}{1.56} & \meanstd{13.81}{5.45} & \meanstd{15.27}{5.85} & \meanstd{25.62}{6.56} & \meanstd{46.88}{9.09} & \meanstd{52.42}{6.87} \\
GPF-plus & \meanstd{27.08}{6.92} & \meanstd{39.31}{4.55} & \meanstd{51.82}{5.06} & \meanstd{12.83}{4.10} & \meanstd{26.16}{2.85} & \meanstd{29.46}{2.64} & \meanstd{34.89}{8.06} & \meanstd{54.47}{4.51} & \meanstd{58.79}{5.19} \\
\gpromptname & \meanstd{21.03}{4.36} & \meanstd{26.98}{3.41} & \meanstd{30.82}{3.43} & \meanstd{18.00}{2.95} & \meanstd{23.71}{2.39} & \meanstd{26.72}{2.49} & \meanstd{45.35}{10.95} & \meanstd{59.93}{5.76} & \meanstd{64.79}{4.02} \\
\midrule\rowcolor{black!5}
\textbf{\method} & \bestmeanstd{53.34}{8.94} & \bestmeanstd{69.38}{2.91} & \bestmeanstd{72.04}{1.77} & \bestmeanstd{44.26}{5.75} & \bestmeanstd{58.19}{4.67} & \bestmeanstd{63.03}{2.12} & \bestmeanstd{57.37}{10.07} & \bestmeanstd{68.21}{4.08} & \bestmeanstd{70.39}{2.47} \\
\bottomrule
\end{tabular*}
\par\vspace{3pt}
\begin{minipage}{\linewidth}\fontsize{8}{9.5}\selectfont\raggedright \lang{Statistics and emphasis follow Table~\ref{tab:main-accuracy}. Each column is a separate dataset/shot condition.}{统计与强调规则同表~\ref{tab:main-accuracy}；每列对应独立的数据集／shot 条件。}\lang{ $\dagger$: GraphPrompt under the common-protocol adaptation.}{ $\dagger$：GraphPrompt 的统一协议适配实现。}\end{minipage}
\end{table*}

\begin{table*}[t]
\centering
\linespread{1}\selectfont
\captionsetup{font=small,skip=5pt}
\caption{\lang{\textbf{Prototype-routing ablation.} Matched prompt recipe and readout.}{\textbf{原型路由消融。}提示训练配方与读出保持一致。}}
\label{tab:ablation}
\fontsize{9}{11}\selectfont
\setlength{\tabcolsep}{1pt}
\renewcommand{\arraystretch}{1.18}
\setlength{\heavyrulewidth}{0.8pt}
\setlength{\lightrulewidth}{0.4pt}
\begin{tabular*}{\linewidth}{@{\extracolsep{\fill}}l*{9}{r}@{}}
\toprule
\lang{Method}{方法} & \multicolumn{3}{c}{\textbf{Cora}} & \multicolumn{3}{c}{\textbf{CiteSeer}} & \multicolumn{3}{c}{\textbf{PubMed}} \\
\cmidrule(lr){2-4}\cmidrule(lr){5-7}\cmidrule(l){8-10}
 & 1-shot & 3-shot & 5-shot & 1-shot & 3-shot & 5-shot & 1-shot & 3-shot & 5-shot \\
\midrule
\multicolumn{10}{@{}l}{\textit{Accuracy (\%) $\uparrow$}} \\[2pt]
\frozenname & \secondmeanstd{24.64}{4.63} & \meanstd{30.86}{3.27} & \meanstd{35.16}{2.40} & \secondmeanstd{21.69}{3.47} & \secondmeanstd{24.96}{3.51} & \secondmeanstd{27.52}{3.72} & \secondmeanstd{53.72}{10.60} & \secondmeanstd{62.49}{9.18} & \secondmeanstd{65.68}{5.56} \\
\sharedname & \meanstd{24.15}{4.41} & \secondmeanstd{31.25}{3.68} & \secondmeanstd{36.30}{2.72} & \meanstd{20.64}{3.07} & \meanstd{23.71}{3.88} & \meanstd{25.60}{3.82} & \meanstd{52.61}{10.70} & \meanstd{60.89}{9.03} & \meanstd{65.00}{5.72} \\
\midrule\rowcolor{black!5}
\textbf{\method} & \bestmeanstd{53.80}{8.45} & \bestmeanstd{69.97}{3.53} & \bestmeanstd{72.87}{3.36} & \bestmeanstd{43.92}{11.01} & \bestmeanstd{59.67}{5.98} & \bestmeanstd{64.70}{3.53} & \bestmeanstd{59.77}{7.47} & \bestmeanstd{65.09}{8.21} & \bestmeanstd{68.34}{6.97} \\
\midrule
\multicolumn{10}{@{}l}{\textit{Macro-F1 (\%) $\uparrow$}} \\[2pt]
\frozenname & \secondmeanstd{23.79}{4.61} & \secondmeanstd{30.11}{2.90} & \meanstd{34.48}{2.30} & \secondmeanstd{20.84}{3.15} & \secondmeanstd{23.85}{3.32} & \secondmeanstd{26.11}{3.34} & \secondmeanstd{52.29}{9.78} & \secondmeanstd{61.43}{9.05} & \secondmeanstd{64.19}{5.48} \\
\sharedname & \meanstd{22.78}{4.10} & \meanstd{30.03}{3.70} & \secondmeanstd{34.69}{2.61} & \meanstd{19.75}{2.74} & \meanstd{22.28}{3.34} & \meanstd{24.42}{3.38} & \meanstd{51.98}{10.37} & \meanstd{59.21}{10.06} & \meanstd{63.59}{6.09} \\
\midrule\rowcolor{black!5}
\textbf{\method} & \bestmeanstd{51.82}{7.47} & \bestmeanstd{68.37}{4.08} & \bestmeanstd{72.02}{2.74} & \bestmeanstd{41.48}{10.31} & \bestmeanstd{56.34}{5.54} & \bestmeanstd{61.34}{3.09} & \bestmeanstd{59.39}{7.93} & \bestmeanstd{64.82}{8.80} & \bestmeanstd{67.80}{7.59} \\
\bottomrule
\end{tabular*}
\par\vspace{3pt}
\begin{minipage}{\linewidth}\fontsize{8}{9.5}\selectfont\raggedright \lang{Local original-ProG runs with support/test overlap: mean $\pm$ sample std. over 20 tasks after averaging three training seeds per task. \textbf{Bold}/\underline{underline}: best/second in each condition. No cross-condition average.}{本地原ProG实现，保留支持／测试重叠：先平均每任务3个训练seed，再报告20个任务的均值 $\pm$ 样本标准差。\textbf{粗体}／\underline{下划线}表示同条件最优／次优；不跨条件平均。}\lang{ Frozen: no prompt; shared prompt: uniform allocation; PRPrompt: prototype routing.}{ 冻结预训练：无提示；共享提示：均匀分配；PRPrompt：原型路由。}\end{minipage}
\end{table*}

本地实验在 Cora、CiteSeer、PubMed 上分别评估每类 1、3、5 个支持节点。官方训练掩码定义支持池，官方测试掩码定义固定查询集。同一次支持抽样使用相同的逐类排列构成嵌套支持集。每个本地条件包含 20 个支持抽样 seed 与 3 个训练 seed；本地方法共享图、预训练权重及支持／查询身份。每个数据集采用原始的边预测预训练权重，其预训练 seed 未记录。

本地对照包含：\textbf{冻结预训练}，不添加提示并采用与 \method 相同的最终原型读出；\textbf{全量微调}，更新预训练编码器及任务读出。提示基线 Accuracy 引用 \citet{mtg2026} 表 4--6 中的 GPPT、Gprompt（GraphPrompt）、All-in-one、GPF、GPF-plus，并保留原表的监督学习与微调参照，排除其 MTG 方法。原文基线结合 ProG 结果与作者复现，报告不同预训练策略中的最佳表现，但未给出各单元对应的具体策略。

\paragraph{文献协议。}
来源论文使用 90\% 测试集、5 次独立任务抽样与随机超参数搜索，主要编码器为两层、隐藏维度 128 的 GCN。这些数值与本地测试划分、权重及训练预算不匹配，因此表~\ref{tab:main-accuracy} 分为文献和本地两个面板。文献均值与标准差原样保留，不重新聚合，也不进行跨面板配对检验或统一排名。原文表格仅提供 Accuracy，其他指标表仅展示本地测量。

\method 与共享提示消融使用同一固定配方。本地全量微调在 3-shot 和 5-shot 下使用三个支持内折，候选为学习率 $\{0.001,0.01\}$ 与权重衰减 $\{0,0.0001\}$ 的四种组合。各候选拟合及最终拟合均为 100 步；1-shot 使用固定参数。因此，本地方法的总优化预算并不完全相同，其余设置见附录~\ref{app:protocol}。

\paragraph{统计单位。}
在固定本地条件内，先对同一支持抽样下的 3 个训练 seed 求均值，再对 20 个支持抽样结果计算均值与样本标准差。不跨数据集或 shot 汇总平均表现。本地配对区间以支持抽样为块进行 10,000 次 bootstrap，报告固定图和权重条件下、未经多重比较校正的 95\% 区间，仅用于冻结预训练、全量微调及共享提示对照。正文报告 Accuracy 与 Macro-F1，其余本地指标见附录~\ref{app:more}。

\subsection{RQ1：完整适配的效果}
本地比较检验相同图、权重及支持标签预算下的适配效果。表~\ref{tab:main-accuracy} 另外加入独立协议的文献参照；表~\ref{tab:main-macro_f1} 报告本地 Macro-F1。每列均为独立条件，不包含跨条件平均行。

Cora 的 5-shot Accuracy 为 73.02\%，冻结预训练与全量微调分别为 33.56\% 和 67.50\%；CiteSeer 在相同 shot 数下分别为 66.17\%、27.38\%、47.89\%。PubMed 的冻结预训练更强：5-shot 为 68.37\%，\method 为 70.98\%，全量微调为 63.96\%。因此，相对冻结预训练的改善在 PubMed 上明显较小，跨数据集平均会掩盖这一依赖关系。

\paragraph{文献报告的提示结果。}
引入文献数据后，不能声称 Accuracy 全面领先。CiteSeer 上 GPF-plus 在 1、3、5-shot 分别报告 59.67\%、72.48\%、75.73\%，本地 \method 对应为 47.12\%、61.20\%、66.17\%。Cora 的 5-shot 本地结果 73.02\% 高于文献 Gprompt 的 69.03\%；PubMed 的 3-shot 本地结果 68.66\% 低于文献 GPF 的 71.20\%。这些属于跨来源的描述性比较，不是匹配协议下的方法效应估计。路由价值由下述本地消融支持，本地配对差值见附录~\ref{app:paired}。

\subsection{RQ2：原型路由的价值}
核心消融将式~\eqref{eq:injection} 的 $B$ 替换为元素均为 $1/C$ 的常数矩阵，保留提示矩阵、尺度、注入位置、优化与最终读出，移除节点相关的提示分配。常数矩阵直接用于组合，不再中心化，否则提示会被完全消除；同时保留无提示的冻结预训练条件。

表~\ref{tab:ablation} 表明，单纯学习共享提示不足以解释完整方法的收益。5-shot 下，原型路由相对共享提示在 Cora、CiteSeer、PubMed 的 Accuracy 分别提高 36.94、40.45、3.02 个百分点，对应配对 95\% 区间为 $[34.87,39.06]$、$[39.06,41.84]$、$[1.39,4.51]$ 个百分点。在各个 1-shot 和 3-shot 条件中，两项主要指标的配对改善也均为正。

该比较支持完整原型路由构造相对共享提示的作用，不能将收益单独归因于原型校正、中心化或图聚合中的某一步。均匀组合虽然保留名义提示参数量，却降低有效自由度，因此也不能视为完全控制容量后对唯一路由机制的证明。

\section{局限性}
实验覆盖三个引文网络，每个图仅使用一份固定预训练权重，任务为同图已知类别分类。预训练 seed 未知，未估计不同预训练运行之间的波动。报告范围在实验完成后确定，因此属于描述性分析，不是新的独立确认实验；结果不支持对图分类、未见类别或任意预训练来源的推广结论。

文献基线使用不同测试划分、5 次抽样以及跨预训练策略选择，其汇总值可提供参照，但不能相对本地固定权重运行隔离适配方法的影响，也不支持配对显著性检验。本地微调选参产生的预算也不同于 \method 的固定配方。路由消融改变了有效提示容量，未区分路由内部各步骤的贡献。此外，方法利用查询特征和全图统计，缺少这类传导式信息的部署条件需另行评估。

\section{结论}
\method 利用支持集锚定的原型组织可学习隐藏层提示的分配，从而适配冻结图编码器。在所评估的 Cora、CiteSeer、PubMed 条件下，逐条件主比较与共享提示消融共同支持这种分配方式的价值。现有证据支持原型路由作为当前协议下的适配设计，其对其他图类型和预训练来源的普适性仍需进一步验证。

\section*{潜在影响}
本文研究公开引文网络基准，不评估涉及个人的实际决策。图聚合与未标注原型校正可能传播错误，尤其当图结构或类别结构偏离所评估设置时。涉及重要后果的应用需要开展基准之外的领域评估。
